% ============================================================
%  part3/ch16-cahn-hilliard.tex  —  Chapter 16: Cahn-Hilliard
% ============================================================

\chapter{Cahn--Hilliard Dynamics}
\label{ch:cahn-hilliard}

\epigraph{%
  Conservation is the source of complication.%
}{---}

\section{Conservation changes everything}

Allen--Cahn allows $\int\phi\,dx$ to change freely.
Many physical quantities must satisfy
$\frac{d}{dt}\int\phi\,dx = 0$:
mass, charge, lamphron density.

\begin{definition}[Cahn--Hilliard equation]
  For conserved $\phi$ with chemical potential
  $\mu = \delta\mathcal{F}/\delta\phi$ and flux
  $\mathbf{J} = -M\nabla\mu$:
  \[
    \boxed{\partial_t\phi
    = M\Delta\mu
    = -M\epsilon^2\Delta^2\phi + M\Delta W'(\phi).}
  \]
\end{definition}

\begin{theorem}[Mass conservation]
  \label{thm:ch-mass}
  \statusproven{}
  \[
    \frac{d}{dt}\int\phi\,dx = 0.
  \]
\end{theorem}

\begin{proof}
  $\frac{d}{dt}\int\phi\,dx = M\int\Delta\mu\,dx = 0$
  by the divergence theorem and Neumann boundary conditions.
\end{proof}

\begin{theorem}[Free-energy dissipation]
  \label{thm:ch-dissipation}
  \statusproven{}
  \[
    \frac{d\mathcal{F}}{dt}
    = -M\int|\nabla\mu|^2\,dx \leq 0.
  \]
\end{theorem}

\begin{proof}
  $\frac{d\mathcal{F}}{dt}
  = \int\mu\,\partial_t\phi\,dx
  = M\int\mu\,\Delta\mu\,dx
  = -M\int|\nabla\mu|^2\,dx \leq 0$.
\end{proof}

\section{Lifshitz--Slyozov scaling}

\begin{theorem}[Lifshitz--Slyozov scaling]
  \label{thm:ls-scaling}
  Conserved phase separation obeys $\ell(t) \sim t^{1/3}$,
  so $z = 3$.
\end{theorem}

\begin{proof}
  Chemical potential scales as $\mu \sim 1/\ell$ (curvature).
  Flux $J \sim |\nabla\mu| \sim 1/\ell^2$.
  Mass transport rate: $d\ell/dt \sim J/\ell \sim 1/\ell^2$.
  Integrating: $\ell^3 \sim t$, so $\ell(t) \sim t^{1/3}$.
\end{proof}

\begin{TechNote}[Conservation changes universality class]
  The conservation law changes $z$ from 2 to 3.
  In RSVP, the lamphron density $L = \ell^2$ satisfies a
  modified conservation equation with advection and source
  terms.
  This is why $z_{\mathrm{RSVP}}$ is expected to differ from
  both pure Allen--Cahn ($z=2$) and pure Cahn--Hilliard ($z=3$).
\end{TechNote}

\section{RSVP corrections to the scaling exponent}

\begin{conjecture}[RSVP coarsening correction]
  \label{conj:rsvp-scaling}
  \statusconjectured{}
  Assume RSVP dynamics modify the coarsening law to
  \[
    \frac{d\ell}{dt} = A\ell^{-1} + B\ell^{-m}.
  \]
  Then:
  \begin{itemize}[leftmargin=2em]
    \item If $m > 1$: Allen--Cahn dominates, $z_{\mathrm{RSVP}} = 2$.
    \item If $m < 1$: RSVP correction dominates,
      $z_{\mathrm{RSVP}} = m+1$.
    \item If $m = 1$: logarithmic correction.
  \end{itemize}
\end{conjecture}

The precise value of $m$ depends on the dimensionless
parameters $Pe$, $\Gamma$, and $\Omega$ and requires
simulation or further analysis.

\WhatSurvived{Cahn--Hilliard coarsening}{%
  Total mass $\int\phi\,dx$\\
  Phase identity at late times\\
  Energy (decreasing)%
}{%
  Mass distribution history\\
  Interface trajectories\\
  Short-scale fluctuations%
}{%
  Total mass exactly.\\
  Late-time phase pattern\\
  partially.%
}{%
  $\ker$: all initial conditions\\
  with same total mass and\\
  same final phase pattern.%
}
